% !TeX program = xelatex
\documentclass[preprint,review,12pt,a4paper]{elsarticle}

% ── Page geometry (preprint mode does not load geometry automatically) ──
\usepackage[top=2.5cm,bottom=2.5cm,left=2.5cm,right=2.5cm]{geometry}

% ── Minimal CJK support for Chinese references in bibliography ──
\usepackage[scheme=plain]{ctex}
\setmainfont{Times New Roman}
\setCJKmainfont{SimSun}[BoldFont=SimHei,ItalicFont=KaiTi]

% ── Figures & Tables ──
\usepackage{graphicx}
\usepackage{caption}
\usepackage{subcaption}
\usepackage{float}
\usepackage{booktabs}
\usepackage{multirow}
\usepackage{array}
\usepackage{longtable}
\usepackage{makecell}
\graphicspath{{figures/}}

% ── Math ──
\usepackage{amsmath}
\usepackage{amssymb}

% ── Hyperlinks ──
\usepackage[unicode=true]{hyperref}
\hypersetup{colorlinks=true,citecolor=blue,linkcolor=black,urlcolor=blue}

% ── Caption font (elsarticle defaults \footnotesize; override for readability) ──
\captionsetup{font=small}

\begin{document}

\begin{frontmatter}

\title{Cross-Location Lane Change Risk Prediction Using Multi-Segment Naturalistic Trajectory Data: A Comparative Study of Four Machine Learning Models with SHAP Interpretability Analysis}

\author{Qiong Fan}
\author{Zhaoqing Shen\corref{cor1}}
\cortext[cor1]{Zhaoqing Shen (1981--), Male, Ph.D., Associate Professor, School of Highway, Chang'an University, Xi'an, China. Research interests: road and traffic engineering.}
\affiliation{organization={School of Highway, Chang'an University}, city={Xi'an}, country={China}}

\begin{abstract}
A fundamental question in lane change risk modeling remains unresolved: can a model trained on one road segment generalize reliably to another with different geometric configuration or speed limit? To address this, we systematically evaluate the cross-location generalization performance of four representative machine learning models---XGBoost, Random Forest (RF), Multilayer Perceptron (MLP), and Long Short-Term Memory (LSTM)---across five real-world expressway scenarios (subgrade segments, monolithic bridge segments, and diverging/merging areas) drawn from the CQSkyEyeX naturalistic trajectory dataset. A multi-indicator risk scoring framework integrating mTTC, THW, PET, and ETTC with entropy-weight-based objective weighting and velocity-adaptive correction was first constructed to assign three-tier risk labels (high-risk: 378, medium-risk: 561, low-risk: 635; 1,574 lane change samples total). Models were evaluated under leave-one-location-out cross-validation (5-fold cross-location CV) alongside random 80/20 split baselines. Results demonstrate that tree-based ensemble models substantially outperform neural networks in cross-location generalization: XGBoost achieves the best cross-location performance (Accuracy = 0.764 $\pm$ 0.078, Weighted F1 = 0.764 $\pm$ 0.079), surpassing RF by approximately 3.2 percentage points, while MLP and LSTM exhibit markedly limited cross-location generalization capability. Road geometry exerts a strong modulating effect: inter-subgrade transferability is highest (Loc2 F1 = 0.875), whereas the merging/diverging area proves the most challenging to generalize (Loc5 F1 = 0.630). SHAP analysis consistently identifies minimum time headway (Time\_Headway\_min) as the most stable cross-location risk discriminant (XGBoost high-risk SHAP contribution: +1.251), while rear-distance features (B\_Dist series) exhibit location-conditional importance. Waterfall decomposition reveals complementary model preferences: XGBoost prioritizes TTC-based transient collision risk, whereas RF emphasizes sustained close-following patterns. These findings provide empirical evidence for model selection and location-adaptive deployment of lane change warning systems, and demonstrate that fusing SHAP explanations from multiple tree-based models yields more complete risk driver identification than any single-model interpretation.
\end{abstract}

\begin{keyword}
Lane change risk assessment \sep Machine learning \sep Cross-location generalization \sep SHAP interpretability \sep Surrogate safety measures \sep Naturalistic driving trajectory
\end{keyword}

\end{frontmatter}

% ── 1 Introduction ──
\section{Introduction}

Lane change (LC) maneuvers account for approximately 4\%--10\% of total traffic collisions~\cite{C01,C02}, and their risk assessment must jointly consider longitudinal collision hazards and lateral multi-vehicle interactions~\cite{C03}. Due to the scarcity of historical crash data, researchers have introduced surrogate safety measures (SSMs) such as TTC~\cite{C05}, PET~\cite{C04}, and SDI~\cite{C06} to quantify potential conflicts. However, conventional SSMs exhibit two fundamental limitations in lane change scenarios: most indicators are designed exclusively for car-following contexts and cannot capture the time-varying nature of lateral interactions~\cite{C03,C07}, and single-dimensional measures are insufficient to characterize the multi-risk-source nature of lane change events~\cite{C01}. To address these limitations, several integrated assessment frameworks have been developed: Park et al.~\cite{C08} proposed the Lane Change Risk Index (LCRI) via fault tree analysis, Wu et al.~\cite{C03} developed the Temporal--Spatial Risk Estimation (TSRE) model, and Xue et al.~\cite{C02} introduced the safety margin concept. In recent years, machine learning (ML) methods have been applied to lane change risk prediction: Feng et al.~\cite{C09} achieved F1 = 84\% in merging scenarios using RF-RFE feature selection, and Wen et al.~\cite{C10} leveraged XGBoost + SHAP to reveal risk factors in merging areas. Deep learning architectures such as DS-LSTM~\cite{C11} and CNN-BiLSTM-Attention~\cite{C12} have demonstrated the potential of temporal modeling for risk recognition. The introduction of SHAP~\cite{C13} and TreeSHAP~\cite{C14} provides a unified framework for model interpretability.

A critical question, however, remains systematically unanswered: \textbf{can a lane change risk model trained on one road segment transfer effectively to another segment with different geometric configuration or speed limit?} Existing studies are predominantly confined to single datasets~\cite{C15}, and although limited cross-dataset evidence suggests site-level heterogeneity in risk factors~\cite{C16}, a systematic multi-location, multi-model comparative assessment is still lacking. The answer to this question directly determines whether lane change warning systems require location-specific model retraining, carrying significant engineering and economic implications.

\subsection{Related Work}

Surrogate safety measures (SSMs) have evolved into a comprehensive taxonomy spanning time-based, distance-based, and deceleration-based indicators since Hayward~\cite{C05} introduced TTC~\cite{C18}. TTC frequently fails under the dynamic multi-vehicle interactions inherent in lane changes~\cite{C01}; Minderhoud and Bovy~\cite{C06} developed extended measures including TET and TIT. In integrated assessment, Park et al.~\cite{C08} constructed the LCRI, Wu et al.~\cite{C03} developed the TSRE model, and Xue et al.~\cite{C02} introduced the safety margin. In risk field theory, Liu and Xiang~\cite{C01} proposed the F-F Diagram, Wang et al.~\cite{C07} advanced the paradigm from deterministic to probabilistic assessment, and Hu et al.~\cite{C29} further introduced a unified field perception model accounting for risk impulse effects. Yu et al.~\cite{C26} integrated spatiotemporal indicators via fault tree analysis to construct a dynamic lane change risk quantification method.

ML methods have demonstrated superior performance over traditional approaches in lane change risk prediction. Qi et al.~\cite{C30} achieved multi-lane merging area lane change recognition and trajectory prediction using K-means++ clustering and CNN-LSTM. Xing et al.~\cite{C31} developed a two-stage lane change risk prediction model incorporating car--truck interaction pressure. Feng et al.~\cite{C09} achieved F1 = 84\% in merging scenarios using RF-RFE feature selection; Wen et al.~\cite{C10} employed XGBoost + SHAP to reveal merging area risk factors. Tree ensemble methods---XGBoost~\cite{C19} and RF~\cite{C20}---are widely adopted due to their strong adaptability to tabular data. Among deep learning approaches, DS-LSTM~\cite{C11} integrates driving style clustering, and CNN-BiLSTM-Attention~\cite{C12} enables multi-vehicle-type lane change risk recognition. Guo et al.~\cite{C28} realized driver distraction detection based on domain adaptation. SHAP~\cite{C13} and TreeSHAP~\cite{C14} provide a theoretical framework for model interpretability, although existing studies predominantly remain at the level of global feature importance ranking without addressing the cross-location dimension.

In the direction of cross-location generalization, Huang et al.~\cite{C22} found that model reliability dropped significantly when deployed from bridge/tunnel segments to merging segments; Wang et al.~\cite{C16} confirmed site-level heterogeneity of risk factors. Xia et al.~\cite{C32} achieved real-time expressway crash risk prediction using abnormal driving behavior data, and Wei et al.~\cite{C34} constructed a deep learning-based traffic conflict prediction method. The domain adaptation strategy of Ahmed et al.~\cite{C23} yielded only limited improvement (3\%--5\%).

Three limitations characterize existing research: (1) most studies rely on simulation environments, lacking coverage of diverse real-world road geometries; (2) a systematic comparison of tree-based and neural network generalization behavior is absent; and (3) SHAP-based explanations have not been leveraged to interpret the internal mechanisms of generalization degradation through the lens of dynamic feature importance variation across locations. To address these gaps, this study conducts a systematic cross-location generalization investigation of lane change risk ML prediction models using five real-world scenarios with distinct road geometries from the CQSkyEyeX dataset~\cite{C17}. The main contributions are:

\begin{enumerate}
  \item \textbf{Multi-location evaluation benchmark}: Based on 1,574 lane change samples spanning three road types (subgrade segments, monolithic bridge segments, and diverging/merging areas), we establish a complete annotation pipeline from multi-indicator weighted safety scoring to three-tier risk labeling (high-risk: 378, medium-risk: 561, low-risk: 635), providing a unified data foundation for cross-location generalization assessment.
  \item \textbf{Four-model cross-location comparison}: Through a paired design of leave-one-location-out cross-validation (5-fold cross-location CV) and random split baselines, we quantify the cross-location performance degradation of four representative models (XGBoost, RF, MLP, LSTM), revealing fundamental differences in generalization behavior between tree ensemble and neural network paradigms.
  \item \textbf{SHAP-driven risk factor universality analysis}: Using multi-model SHAP consistency analysis across XGBoost and RF, we separate universal risk factors (generalizable across locations) from conditional factors (location-specific) for the first time in the cross-location dimension, providing feature-level evidence for location-adaptive model design.
\end{enumerate}

% ── 2 Methods ──
\section{Lane Change Risk Prediction Model Development}

\subsection{CQSkyEyeX Naturalistic Trajectory Data: Source and Preprocessing}

This study employs the CQSkyEyeX naturalistic driving trajectory dataset~\cite{C17}\footnote{Dataset portal: http://www.cqskyeyex.com/CQSkyEyeX}. The dataset extracts vehicle trajectories from UAV aerial videos via computer vision methods~\cite{C24}, with positioning accuracy RMSE $<0.10$ m. The original sampling rate is 30 Hz; all data were uniformly down-sampled to 5 Hz for this study.

We selected the first five publicly annotated and released scenarios from the dataset, whose road geometries span three typical expressway segment types: subgrade (basic) segments, monolithic bridge segments, and diverging/merging areas. Table~\ref{tab:scenes} summarizes the detailed characteristics of each scenario.

\begin{table}[H]
\centering
\caption{Characteristics of the five CQSkyEyeX scenarios}
\label{tab:scenes}
\small
\begin{tabular}{lcccccc}
\toprule
Scenario & Segment type & Lanes (one-way) & Speed limit (km/h) & Segment length (m) & LC samples \\
\midrule
Loc1 & Subgrade basic & \multirow{5}{*}{3} & \multirow{4}{*}{120} & \multirow{4}{*}{$\!\sim\!420$} & 151 \\
Loc2 & Subgrade basic &                    &                      &                               & 243 \\
Loc3 & Monolithic bridge &                  &                      &                               & 440 \\
Loc4 & Monolithic bridge &                  &                      &                               & 306 \\
\cmidrule(lr){4-5}
Loc5 & Diverge/merge area &                & 100                  & $\!\sim\!350$                 & 434 \\
\bottomrule
\end{tabular}
\end{table}

Figure~\ref{fig:velocity_distribution} displays the velocity distribution densities for three representative scenarios. Loc1 (subgrade) exhibits a pronounced bimodal structure (a low-speed truck peak and a high-speed passenger-car peak); Loc3 (bridge) shows weakened bimodality; Loc5 (diverging/merging area) tends toward a unimodal concentration. Wang et al.~\cite{C27} similarly documented the differential impact of segment geometry on operating speed distributions, corroborating this observation. This distributional divergence determines the sensitivity of the correction factor $K$ to the exponent $B$: in Loc5, the velocity concentration at lower speeds yields $K \equiv 1$, whereas in Loc1 the high-speed peak group receives effective correction from $K$.

\begin{figure}[H]
\centering
\begin{subfigure}{0.32\textwidth}
\centering
\includegraphics[width=\textwidth]{velocity_following_distribution_location1.png}
\caption{Loc1 (subgrade)}
\end{subfigure}
\begin{subfigure}{0.32\textwidth}
\centering
\includegraphics[width=\textwidth]{velocity_following_distribution_location3.png}
\caption{Loc3 (bridge)}
\end{subfigure}
\begin{subfigure}{0.32\textwidth}
\centering
\includegraphics[width=\textwidth]{velocity_following_distribution_location5.png}
\caption{Loc5 (diverge/merge area)}
\end{subfigure}
\caption{Velocity distribution density curves of following vehicles in three representative scenarios. Loc1 exhibits a clear bimodal pattern (truck low-speed peak vs.\ passenger-car high-speed peak); Loc3 shows weakened bimodality; Loc5 tends toward unimodal concentration.}
\label{fig:velocity_distribution}
\end{figure}

The original dataset provides neighboring vehicle IDs; spatial distances to neighboring vehicles in each direction are computed from coordinate geometry. Lane change samples are extracted using a reference window spanning $t_{-100}$ to $t_{-25}$ (75 frames, 3~s~\cite{C02,C03}) centered on the lane-crossing moment $t_0$. After moving-average smoothing and denoising, 1,574 valid samples were obtained. The temporal non-stationarity of risk indicators within each lane change window motivates a feature engineering strategy employing aggregate statistics (mean, standard deviation, extrema) to capture extreme-state information rather than single-frame instantaneous values.

\subsection{Entropy-Weight-Based Surrogate Safety Indicators and Risk Scoring Framework}

To achieve unsupervised risk annotation of lane change trajectories while overcoming the limitations of traditional discrete-threshold scoring (brittle class boundaries, insensitivity to small indicator variations), we adopt a continuous risk scoring scheme based on exponential decay functions. The weights $\{w_i\}$ in the scoring formula are objectively determined from the data via the Entropy Weight Method (EWM). EWM is a label-free objective weighting approach whose core logic is: indicators with more dispersed information (greater inter-vehicle variation in activation values) possess higher discriminative power and receive higher weights. The computation proceeds as follows:

\textbf{Step 1 --- Risk activation matrix.} For $n$ lane-changing vehicles and $m = 5$ SSM indicators (mTTC, THW, PET, F\_ETTC, OL\_PET), the exponential decay values of the most dangerous frame within the full observation window are taken for each vehicle to construct $\mathbf{X} \in [0,1]^{n \times m}$:
\begin{equation}
X_{ij} = \max_{t \in [1,\,75]}\; \exp\!\left(-\frac{v_j^{(i)}(t)}{k_j}\right)
\end{equation}
If the $j$-th indicator has no corresponding neighboring vehicle for the $i$-th vehicle, then $X_{ij} = 0$.

\textbf{Step 2 --- Proportion matrix.} Column-normalize to obtain $\mathbf{P}$, where $P_{ij} = X_{ij} / \sum_{i=1}^n X_{ij}$:
\begin{equation}
P_{ij} = \frac{X_{ij}}{\sum_{i=1}^{n} X_{ij}}
\end{equation}

\textbf{Step 3 --- Information entropy and divergence coefficient.} The information entropy of the $j$-th indicator is $E_j = -\frac{1}{\ln n}\sum_i P_{ij}\ln P_{ij}$, with divergence coefficient $D_j = 1 - E_j$. Larger $E_j$ (more concentrated information) yields lower weight.

\textbf{Step 4 --- Weight normalization.} $w_j = D_j / \sum_k D_k$, satisfying $\sum w_j = 1$.

The EWM weights depend solely on the risk activation matrix $\mathbf{X}$ of the five SSM indicators and have no access to the 60-dimensional sample features or risk labels used for model training, guaranteeing the statistical independence and reproducibility of weight calibration. The resulting lane change vehicle weights are $\{w_j\}^{5}_{j=1} = \{0.247, 0.083, 0.131, 0.356, 0.183\}$ (mTTC/THW/PET/F\_ETTC/OL\_PET), and following vehicle weights are $\{0.463, 0.114, 0.423\}$ (mTTC/THW/B\_mTTC). Notably, F\_ETTC receives the highest weight (0.356), substantially exceeding its empirically estimated value (0.20), indicating that forward extended collision time exhibits the greatest inter-vehicle discriminative power---i.e., the heterogeneity of forward collision risk across different lane change events far surpasses that of the other four indicators. Conversely, THW receives only 0.083, reflecting that time headway distributions are highly concentrated across vehicles ($E_j$ close to 1) and provide limited additional discriminative information.

\subsubsection{Vehicle Interaction Topology in Lane Change Scenarios}

During a lane change maneuver, the subject vehicle (SV) and four surrounding vehicles constitute a typical five-vehicle interaction system (Fig.~\ref{fig:lc_interaction}): the preceding vehicle in the current lane (PV), the following vehicle in the current lane (FV), the lead vehicle in the target lane (LV), and the lag vehicle in the target lane (LgV). Different SSM indicators quantify the conflict urgency between SV and specific neighboring vehicles; clarifying this mapping is the physical prerequisite for the subsequent risk scoring formulation.

\begin{figure}[H]
\centering
\includegraphics[width=0.6\textwidth]{fig_lc_interaction.png}
\caption{Interaction topology of the subject vehicle (SV) with four surrounding vehicles. LV: Lead Vehicle (target lane); LgV: Lag Vehicle (target lane); PV: Preceding Vehicle (current lane); FV: Following Vehicle (current lane). Arrows indicate travel direction and lateral displacement intent.}
\label{fig:lc_interaction}
\end{figure}

The correspondence between the five SSM indicators and the vehicle pairs defined above is as follows: mTTC and THW quantify the longitudinal collision risk between SV and PV (rear-end risk in the current lane); PET measures the temporal proximity with which SV and a target-lane neighbor (LV or LgV) successively pass the same spatial point; F\_ETTC extends collision time to the two-dimensional plane, measuring the combined approach rate between SV and LV along the collision hazard direction; OL\_PET quantifies the lateral time margin between SV's encroachment into the target lane and the arrival of LgV at the same crossing point. In the following-vehicle risk score (Eq.~8), B\_mTTC measures the rear-end collision risk between FV and SV.

\subsubsection{Mathematical Definitions of Surrogate Safety Indicators}

\textbf{(1) Modified Time-to-Collision (mTTC).} Conventional TTC assumes constant motion and neglects acceleration; mTTC incorporates acceleration into collision time estimation. Let the relative position vector be $\Delta\mathbf{s} = (\Delta x, \Delta y)$, the relative velocity vector be $\Delta\mathbf{v} = (\Delta v_x, \Delta v_y)$, and the relative acceleration vector be $\Delta\mathbf{a} = (a_x, a_y)$. Then:

\begin{equation}
\text{mTTC} = \frac{-\langle\Delta\mathbf{v}, \Delta\mathbf{s}\rangle + \sqrt{\langle\Delta\mathbf{v}, \Delta\mathbf{s}\rangle^2 + 2\langle\Delta\mathbf{a}, \Delta\mathbf{s}\rangle\|\Delta\mathbf{s}\|^2}}{2\langle\Delta\mathbf{a}, \Delta\mathbf{s}\rangle}
\end{equation}

where $\langle\cdot,\cdot\rangle$ denotes the inner product of vectors. When the following vehicle's velocity does not exceed that of the preceding vehicle (zero denominator or negative radicand), $\text{mTTC} \to \infty$, indicating that collision is impossible under the current motion state.

\textbf{(2) Time Headway (THW).} Characterizes the time required for the subject vehicle to reach the current position of the preceding vehicle at its current speed:

\begin{equation}
\text{THW} = \frac{d_{\!f}}{v_e}
\end{equation}

where $d_{\!f}$ is the longitudinal distance from the front of the subject vehicle to the rear of the preceding vehicle (m), and $v_e$ is the subject vehicle speed (m/s). THW directly reflects the driver's available reaction time reserve.

\textbf{(3) Post-Encroachment Time (PET).} Measures the time difference between two vehicles successively passing the same conflict point in space~\cite{C04}:

\begin{equation}
\text{PET} = t_{2} - t_{1}
\end{equation}

where $t_1$ is the moment the first vehicle clears the conflict zone, and $t_2$ is the moment the second vehicle arrives at the conflict zone. Smaller PET indicates more urgent conflict.

\textbf{(4) Forward Extended Time-to-Collision (F\_ETTC).} Extends collision time from the longitudinal dimension to the two-dimensional plane~\cite{C03}. Let the subject vehicle position be $\mathbf{s}_e$, the preceding vehicle position be $\mathbf{s}_f$, their relative velocity vector be $\mathbf{v}_{ef} = \mathbf{v}_e - \mathbf{v}_f$, and the angle $\theta = \angle(\mathbf{v}_{ef},\; \mathbf{s}_f - \mathbf{s}_e)$. Then:

\begin{equation}
\text{F\_ETTC} = \frac{\|\mathbf{s}_f - \mathbf{s}_e\|}{\|\mathbf{v}_{ef}\|} \cdot \cos\theta
\end{equation}

The $\cos\theta$ term excludes the component of the lateral relative velocity that is not aligned with the SV--LV connecting direction, making F\_ETTC equivalent to the collision time projected onto the collision hazard direction.

\textbf{(5) Over-Lane Post-Encroachment Time (OL\_PET).} For lane change scenarios involving an adjacent lane, measures the time difference between the moment the lane-changing vehicle (SV) crosses the lane marking and the moment the following vehicle (FV) in the original lane arrives at the same crossing point:

\begin{equation}
\text{OL\_PET} = t_{\text{FV}} - t_{\text{LC}}
\end{equation}

where $t_{\text{LC}}$ is the moment the leading edge of the lane-changing vehicle first contacts the lane marking boundary, and $t_{\text{FV}}$ is the moment the leading edge of the following vehicle in the same lane arrives at that spatial point. Smaller OL\_PET indicates that the lane change maneuver poses more urgent disruption to normally traveling vehicles in the original lane.

\subsubsection{Lane Change Risk Scoring Function}

For a single lane change sample spanning a 75-frame temporal sequence, let the observed value of indicator $i$ at time $t$ be $v_i(t)$. Each indicator value is mapped to a risk contribution within $[0, w_i]$:

\begin{equation}
S_{\text{LC}}(t) = K \cdot \sum_{i \in \mathcal{M}} w_i \cdot \exp\!\left(-\frac{v_i(t)}{k_i}\right)
\end{equation}

\begin{equation}
S_{\text{LC}} = K \cdot \sum_{i \in \mathcal{M}} w_i \cdot \max_{t \in [1,\,75]}\;\exp\!\left(-\frac{v_i(t)}{k_i}\right)
\end{equation}

Equation~(6) is the frame-level risk score and Eq.~(7) is the final per-vehicle risk score (taking the maximum contribution of each indicator from the most dangerous frame across the full window, then summing). $\mathcal{M} = \{\text{mTTC}, \text{THW}, \text{PET}, \text{F\_ETTC}, \text{OL\_PET}\}$, with EWM objective weights $\{w_i\} = \{0.247, 0.083, 0.131, 0.356, 0.183\}$ (standard EWM output, not tuned). $K$ is the velocity correction factor defined in the following subsection.

\subsubsection{Calibration of Decay Constants $k_i$}

The decay constant $k_i$ controls the sensitivity of $\exp(-v/k)$ to indicator variation. Collision-time-type indicators (mTTC, PET, F\_ETTC, OL\_PET) adopt $k = 12$~s, whose ratio relative to the TTC danger threshold of 2--5~s~\cite{C06} ensures that contributions vary smoothly within the 0.3--0.9 range. THW adopts $k = 6$~s because its observation interval (0--8~s) is narrower and risk escalates sharply when THW $< 1$~s~\cite{C18}.

\subsubsection{Following-Vehicle Risk Scoring}

For following vehicles (preceding or following) interacting with the target lane-changing vehicle, risk scoring employs a reduced indicator set and different weight distribution, reflecting the fact that car-following scenarios involve only longitudinal collision risk:

\begin{equation}
S_{\text{FW}} = K \cdot \sum_{i \in \mathcal{M}_{\text{FW}}} w_i \cdot \max_{t \in [1,\,75]}\;\exp\!\left(-\frac{v_i(t)}{k_i}\right)
\end{equation}

where $\mathcal{M}_{\text{FW}} = \{\text{mTTC}, \text{THW}, \text{B\_mTTC}\}$, with EWM weights $\{w_i\} = \{0.463, 0.114, 0.423\}$. B\_mTTC is the backward modified time-to-collision (mTTC with respect to the following vehicle), sharing the same decay constant $k = 12.0$~s. The car-following weights tilt heavily toward mTTC and B\_mTTC (combined 0.886), since front/rear collision is the predominant crash mode in car-following scenarios, concentrating the information gain. The car-following score shares the same velocity correction factor $K$ as the lane change score, enabling direct comparison between the two.

\subsubsection{Velocity Correction Factor $K$ and Parameter $B = 1.3$}

Different road speed limits systematically shift the baseline operating speed of vehicles; identical SSM values at different speed limits do not correspond to equivalent risk levels. For example, THW = 2~s at a 120~km/h limit ($v_e \approx 33.3$~m/s, spacing $\approx 67$~m) versus at a 100~km/h limit ($v_e \approx 27.8$~m/s, spacing $\approx 56$~m) carries different physical safety margins. To address this, a velocity-adaptive nonlinear correction factor is introduced:

\begin{equation}
K = \begin{cases}
1, & v_{85} \leq v_0 \\[4pt]
\displaystyle\left(\frac{v_{85}}{v_0}\right)^{B}, & v_{85} > v_0
\end{cases}
\end{equation}

where $v_{85}$ is the 85th percentile of the 75-frame velocity sequence of a single vehicle (m/s), $v_0$ is the posted speed limit of the road segment (m/s), and the exponent $B = 1.3$.

The choice $B = 1.3$ represents a compromise between linear correction ($B = 1$, where only the SSM itself carries speed information) and kinetic-energy-squared correction ($B = 2$, where $K \propto v^2$). Pure squared correction would doubly amplify scores in high-speed scenarios (since time-based indicators already embed speed information), while zero correction would cause aggressive and normal driving behaviors in different speed-limit scenarios to converge in score. The $(v_{85}/v_0)^{1.3}$ form ensures that correction is negligible near the speed limit ($K \approx 1$) with substantive penalty only under pronounced speeding.

\subsubsection{Three-Tier Risk Label Mapping}

The continuous score $S_{\text{LC}}$ (or $S_{\text{FW}}$) is mapped to three discrete tiers suitable for supervised learning via fixed thresholds:

\begin{table}[H]
\centering
\caption{Risk-tier thresholds for different scenarios}
\label{tab:risk_thresholds}
\small
\begin{tabular}{lcc}
\toprule
Scenario & Medium-risk threshold & High-risk threshold \\
\midrule
Lane change ($S_{\text{LC}}$) & 0.40 & 0.60 \\
Following ($S_{\text{FW}}$)   & 0.20 & 0.35 \\
\bottomrule
\end{tabular}
\end{table}

The thresholds for car-following scenarios are lower than those for lane change scenarios because the car-following indicator set contains only three indicators, all directly collision-related, resulting in a systematically lower score baseline than the five-indicator lane change score. Adopting scenario-specific thresholds can be viewed as calibrating for different crash prior probabilities, producing comparable label distributions under equivalent risk levels.

\subsection{Multi-Dimensional Statistical Feature Aggregation from Temporal Trajectories}

For each vehicle's 75-frame temporal trajectory, a statistical aggregation strategy maps the variable-length temporal sequence to a fixed-dimensional feature vector. For the 15 base kinematic indicators (Table~\ref{tab:features}), four aggregate statistics are computed: mean, standard deviation (std), minimum (min), and maximum (max), yielding a total of $15 \times 4 = 60$ features (including the newly added longitudinal jerk, Long\_Jerk, and lateral jerk, Lat\_Jerk). Jerk, the time derivative of acceleration, measures driving smoothness: high-amplitude jerk fluctuations reflect non-stationary maneuvers such as emergency braking or sudden steering, which are closely related to lane change risk. The mean captures overall behavioral tendency, the standard deviation characterizes driving control stability, and the extrema capture information from the most dangerous moments. Any inf or missing values are uniformly filled with 0.

\begin{table}[H]
\centering
\caption{Base kinematic feature list}
\label{tab:features}
\small
\begin{tabular}{llp{6cm}}
\toprule
Category & Feature & Physical meaning \\
\midrule
Velocity & Velocity, long\_Vel, lat\_Vel & Resultant speed, longitudinal/lateral components \\
Acceleration & long\_Acc, lat\_Acc & Longitudinal/lateral acceleration \\
Distance & \makecell[l]{Following\_dist, B\_Dist, LB\_Dist, \\ RB\_Dist, LF\_Dist, RF\_Dist} & Distance to neighboring vehicles in each direction \\
Jerk & Long\_Jerk, Lat\_Jerk & Longitudinal/lateral jerk (m/s\textsuperscript{3}) \\
Time & TTC, Time\_Headway & Collision time and time headway \\
\bottomrule
\end{tabular}
\end{table}

\subsection{Multi-Paradigm Machine Learning Model Selection}

Four representative models spanning two major paradigms---tree ensembles and neural networks---are selected. XGBoost~\cite{C19} and Random Forest~\cite{C20} represent tree ensemble methods; MLP and LSTM~\cite{C25} represent neural network methods. LSTM directly accepts the raw 75-frame temporal sequence as input, whereas the other three models use the 60-dimensional aggregate feature vector. Detailed hyperparameter configurations are provided in the Appendix.

\subsection{Cross-Location Generalization Evaluation Protocol}

Two complementary evaluation protocols are designed to disentangle a model's intrinsic fitting capacity from its cross-location generalization ability:

\begin{enumerate}
  \item \textbf{Leave-One-Location-Out Cross-Validation (Cross-Location CV)}. All samples from one location are sequentially held out as the test set, with the remaining four locations' samples pooled as the training set, repeated five times such that each location serves as the test set exactly once. This protocol directly simulates the most common real-world deployment scenario for lane change warning systems: when a new road segment goes online, only historical labeled data from existing segments are available for model training. Evaluation metrics---Accuracy, Weighted F1, and Macro F1---are reported as mean $\pm$ standard deviation across the five folds.

  \item \textbf{Random 80/20 Split.} All 1,574 samples are partitioned into 80\% training (1,259 vehicles) and 20\% test (315 vehicles) via stratified random sampling (stratified by risk class). This protocol eliminates inter-location data distribution shift and represents the upper-bound performance achievable within a homogeneous data distribution, serving as the reference baseline for cross-location generalization. The performance gap between the two evaluation protocols directly quantifies the degree of model degradation caused by location-specific feature distribution differences.
\end{enumerate}

In the interpretability dimension, SHAP values are computed for both tree-based models (XGBoost and RF) on all 1,574 samples, yielding beeswarm summary plots and waterfall decomposition plots. Cross-model consistency of feature attributions is examined to identify robust risk factors; differential analysis across individual cross-location folds further reveals location-dependent variations in feature importance.

% ── 3 Results ──
\section{Model Performance Analysis and Validation}

\subsection{Risk Distribution and Spatial Characteristics}

Figure~\ref{fig:risk_by_location} presents the ridgeline density plot of continuous risk distributions for lane-changing (LC) and following (FW) vehicles across the five scenarios. In Loc1--Loc4, the following-vehicle density peak concentrates in the 0--0.4 low-risk interval, while the lane-change peak shifts rightward to the 0.4--0.7 medium-to-high risk interval, indicating that lane change maneuvers systematically incur higher instantaneous collision risk. Loc5 (diverging/merging area) exhibits a fundamental inversion of the distribution pattern: the following-vehicle peak shifts substantially rightward, the lane-change peak concentrates in the 0.6--0.8 high-risk interval, and the two curves overlap extensively in the medium-to-high risk region, suggesting that both following and lane-changing vehicles operate under elevated collision risk in this scenario. A monotonic rightward shift is observed from Loc2 (lowest risk) to Loc5 (highest risk), reflecting the incremental effect of road environment complexity on driving risk.

\begin{figure}[H]
\centering
\includegraphics[width=0.82\textwidth]{risk_by_location.png}
\caption{Ridgeline density plot of lane change and following risk distributions across five scenarios. Light blue: following; brick red: lane change. Loc1--Loc4 lane change curves are right-shifted relative to following; Loc5 shows both types concentrated in the medium--high risk zone. Vertical axis (top to bottom): Loc5 $\to$ Loc1.}
\label{fig:risk_by_location}
\end{figure}

The spatial heatmap of high-risk lane change points (Fig.~\ref{fig:risk_heatmap}) reveals a noteworthy phenomenon: high-risk frames are not uniformly distributed across the lane surface but exhibit pronounced clustering at specific spatial locations---such as near the nose of merging area terminals, around lane crossing points, and in bridge segments with longitudinal gradient transitions. This finding aligns with the observation of Wu et al.~\cite{C03} that ``lane change risk exhibits strong spatial non-uniformity,'' providing direct graphical evidence for the necessity of location-specific risk assessment strategies.

\begin{figure}[H]
\centering
\includegraphics[width=\textwidth]{fig_risk_heatmap.png}
\caption{Spatial heatmap of high-risk lane change points across five scenarios (lane markings delineated in white/dark gray).}
\label{fig:risk_heatmap}
\end{figure}

\subsubsection{Comparison of Risk Levels: Lane Change vs.\ Car-Following}

To verify whether lane change maneuvers are inherently associated with higher instantaneous collision risk than car-following driving, we conducted a head-to-head comparison of the continuous risk scores of lane-changing vehicles and contemporaneous same-segment following vehicles across all five locations (Fig.~\ref{fig:lc_vs_following}), with sample sizes of 1,521 (LC) and 6,457 (FW), respectively.

\begin{figure}[H]
\centering
\includegraphics[width=0.82\textwidth]{risk_score_distribution_box.png}
\caption{Comparison of continuous risk score distributions: lane change vs.\ car-following vehicles. (Left) Distribution histogram excluding zero-score samples; inset bar: proportion of score = 0. (Right) Box plot with diamond mean markers; key statistics annotated in upper right.}
\label{fig:lc_vs_following}
\end{figure}

From the distribution histogram (left panel), the proportion of exactly-zero risk scores among following vehicles (inset bar, 12.4\%) is approximately ten times that among lane-changing vehicles (1.3\%), indicating that a large fraction of car-following scenarios are in a completely conflict-free safe state in actual operation, whereas lane-changing vehicles---due to lateral crossing and interactions with multiple neighboring vehicles---almost never experience zero-risk conditions. After excluding zero-score samples, the density peak interval of non-zero risk scores for lane-changing vehicles (0.3--0.5) is substantially higher than that for following vehicles (0.1--0.2): even after removing conflict-free samples, the conditional risk distribution of lane changes remains globally right-shifted.

The box plot (right panel) reveals differences in distribution shape: the median lane change risk score (0.465) and upper quartile are systematically higher than those for following (0.236), and the whiskers of lane change scores extend to larger extreme values. The proportion of high-risk samples among lane-changing vehicles is 25.0\% (380/1,521), compared to only 11.0\% (712/6,457) among following vehicles---a ratio of approximately 2.3$\times$. Combining the evidence from both dimensions, a robust conclusion emerges: \textbf{under the same traffic environment and unified index-based scoring framework, lane change maneuvers carry systematically higher instantaneous collision risk than car-following driving, a finding that holds consistently across all five locations.}

\subsection{Cross-Location Generalization: Evaluation and Enhancement Strategy Comparison}

To systematically assess the cross-location generalization of lane change risk models, leave-one-location-out cross-validation was first applied to evaluate the baseline performance of the four models (XGBoost/RF/MLP/LSTM, default hyperparameters). Results show that tree ensemble models outperform neural networks, but generalization to the merging/diverging area degrades substantially. To address this limitation, four enhancement strategies were explored: SMOTE oversampling, Optuna hyperparameter optimization, SMOTE+Optuna joint optimization, and PCA dimensionality reduction. Table~\ref{tab:enhance_summary} compares the enhancement strategies: SMOTE yielded a modest improvement in cross-location performance (F1 = 0.741 vs.\ baseline 0.734); Optuna tuning yielded further gains; PCA dimensionality reduction exacerbated degradation; and ST joint optimization achieved the best cross-location F1 (0.764). The following presents detailed results with ST joint optimization as the final configuration.

\begin{table}[H]
\centering
\caption{Enhancement strategy performance comparison (XGBoost)}
\label{tab:enhance_summary}
\small
\begin{tabular}{lcccc}
\toprule
Configuration & Locations & Cross-Loc F1 & Random F1 & Degradation \\
\midrule
Baseline & 5 loc & 0.734 & 0.761 & $-0.027$ \\
SMOTE & 5 loc & 0.741 & 0.794 & $-0.053$ \\
Optuna & 5 loc & 0.762 & 0.780 & $-0.018$ \\
ST Joint & 5 loc & 0.764 & 0.793 & $-0.029$ \\
ST$^*$ (uniform upper bound) & loc1--4 & 0.787 & 0.852 & $-0.065$ \\
\bottomrule
\end{tabular}
\end{table}
\noindent$^*$ST$^*$ is the control experiment excluding Loc5, reflecting the optimization ceiling under homogeneous scenarios.

SMOTE-synthesized samples depend on the internal structure of the training locations and cannot represent the risk characteristics of unseen locations, hence cross-location performance degrades rather than improves; all SMOTE experiments were performed independently within each fold's training set, precluding data leakage. Optuna Bayesian optimization yielded substantial positive effects for tree models (XGBoost Cross-Loc F1 improved from baseline 0.734 to 0.762, a gain of 0.028), but LSTM remained at a low level. ST joint optimization achieved the best cross-location F1 (0.764), only marginally better than Optuna alone, indicating that sample synthesis and hyperparameter search do not exhibit positive synergy for cross-location generalization. After PCA dimensionality reduction (15 $\to$ 12 dimensions), LSTM F1 dropped from 0.598 to 0.509 (Fig.~\ref{fig:pca_variance}), further confirming that temporal correlation patterns carry critical discriminative information, and dimensionality reduction amplifies the inter-location performance gap.

\begin{figure}[H]
\centering
\includegraphics[width=0.6\textwidth]{fig_pca_variance.png}
\caption{PCA cumulative variance curve (15 dimensions $\to$ 12 dimensions, 97.8\% cumulative variance).}
\label{fig:pca_variance}
\end{figure}

Based on the above comparison, ST joint optimization was selected as the optimal configuration. Table~\ref{tab:cross_loc} reports its cross-location generalization performance.

\begin{table}[H]
\centering
\caption{Cross-location CV performance with ST optimization}
\label{tab:cross_loc}
\small
\begin{tabular}{lccc}
\toprule
Model & Accuracy & Weighted F1 & Macro F1 \\
\midrule
\textbf{XGBoost} & \textbf{0.764$\pm$0.078} & \textbf{0.764$\pm$0.079} & \textbf{0.747$\pm$0.063} \\
RF  & 0.730$\pm$0.086 & 0.732$\pm$0.086 & 0.706$\pm$0.079 \\
MLP & 0.659$\pm$0.093 & 0.660$\pm$0.094 & 0.627$\pm$0.077 \\
LSTM & 0.481$\pm$0.106 & 0.463$\pm$0.113 & 0.430$\pm$0.095 \\
\bottomrule
\end{tabular}
\end{table}

Location-specific fold performance varies substantially (Table~\ref{tab:folds}): Loc2 (subgrade) achieves the best generalization (XGBoost F1 = 0.875), while Loc5 (diverging/merging area) performs worst (F1 = 0.630).

\begin{table}[H]
\centering
\caption{Per-fold XGBoost performance (ST optimization)}
\label{tab:folds}
\small
\begin{tabular}{lccccc}
\toprule
Test scenario & Samples & Acc & F1 & MacroF1 \\
\midrule
Loc1 (subgrade)    & 151 & 0.795 & 0.797 & 0.787 \\
Loc2 (subgrade)    & 243 & 0.872 & 0.875 & 0.828 \\
Loc3 (bridge)      & 440 & 0.768 & 0.768 & 0.747 \\
Loc4 (bridge)      & 306 & 0.755 & 0.751 & 0.733 \\
Loc5 (diverge/merge) & 434 & 0.631 & 0.630 & 0.639 \\
\bottomrule
\end{tabular}
\end{table}

\begin{figure}[H]
\centering
\includegraphics[width=0.80\textwidth]{fig_cross_location_comparison.png}
\caption{F1 comparison across models and locations.}
\label{fig:cross_location}
\end{figure}

\begin{figure}[H]
\centering
\begin{subfigure}{0.45\textwidth}
\centering\includegraphics[width=\textwidth]{09_exp_ST_cross_location_radar.png}
\caption{Cross-location CV}
\end{subfigure}
\hfill
\begin{subfigure}{0.45\textwidth}
\centering\includegraphics[width=\textwidth]{09_exp_ST_random_split_radar.png}
\caption{Random 80/20 split}
\end{subfigure}
\caption{Four-model performance radar charts under ST joint optimization.}
\label{fig:radar}
\end{figure}

Under the random 80/20 split (within-distribution evaluation, Table~\ref{tab:random}), XGBoost F1 rises to 0.793, RF to 0.763, MLP to 0.659, and LSTM to 0.569. XGBoost's cross-location degradation of approximately 2.9 percentage points (0.793 $\to$ 0.764) indicates favorable cross-location transferability of its feature--risk mapping.

\begin{table}[H]
\centering
\caption{Random 80/20 split results}
\label{tab:random}
\small
\begin{tabular}{lccc}
\toprule
Model & Accuracy & Weighted F1 & Macro F1 \\
\midrule
\textbf{XGBoost} & \textbf{0.790} & \textbf{0.793} & \textbf{0.782} \\
RF  & 0.759 & 0.763 & 0.750 \\
MLP & 0.654 & 0.659 & 0.647 \\
LSTM & 0.581 & 0.569 & 0.558 \\
\bottomrule
\end{tabular}
\end{table}

The inter-fold variance reflects cross-location performance uncertainty: fold-to-fold fluctuations in Loc3 and Loc5 are markedly higher than those in Loc1 and Loc2, further corroborating that road type is a key factor influencing prediction stability.

\subsection{SHAP-Based Cross-Location Stability Analysis}

\subsubsection{Global Feature Importance}

Table~\ref{tab:shap} summarizes the SHAP contributions of individual features to the three risk classes for the XGBoost model.

\begin{table}[H]
\centering
\caption{XGBoost SHAP Top-5 feature contributions (ranked by overall influence)}
\label{tab:shap}
\small
\begin{tabular}{lcccc}
\toprule
Rank & Feature & High-risk contrib. & Medium-risk contrib. & Low-risk contrib. \\
\midrule
1 & Time\_Headway\_min & +1.251 & +0.309 & +0.481 \\
2 & B\_Dist\_mean      & +0.160 & +0.096 & +0.347 \\
3 & Time\_Headway\_max & +0.037 & +0.048 & +0.369 \\
4 & B\_Dist\_max       & +0.135 & +0.051 & +0.262 \\
5 & Following\_dist\_min & +0.155 & +0.069 & +0.138 \\
\bottomrule
\end{tabular}
\end{table}

The SHAP ranking of Random Forest (Top-5: Time\_Headway\_min, Time\_Headway\_mean, Following\_dist\_min, Following\_dist\_max, Time\_Headway\_max) is highly consistent with that of XGBoost, indicating that the following core findings exhibit cross-model robustness:

\begin{enumerate}
  \item \textbf{Time\_Headway-series features dominate absolutely.} In XGBoost, the high-risk contribution of Time\_Headway\_min (+1.251) far exceeds that of the second-ranked feature (+0.160) by more than an order of magnitude. This result is mechanistically interpretable: minimum time headway directly quantifies the shortest reaction time available between the subject vehicle and the preceding vehicle during a lane change---once this headway falls below the driver perception--reaction threshold (typically 1.0--2.0~s), collision becomes nearly unavoidable. Wen et al.~\cite{C10} similarly identified mean time headway as an important feature in merging area risk discrimination; the present study further confirms that the minimum (rather than the mean) is the risk indicator variant with the strongest discriminative power.

  \item \textbf{Rear-distance (B\_Dist) features constitute the second tier.} In the XGBoost ranking, B\_Dist\_mean and B\_Dist\_max rank 2nd and 4th, respectively. Rear distance characterizes the spatial separation between the subject vehicle and the following vehicle during a lane change: smaller distances imply higher potential severity of rear-end collision should the following vehicle fail to adjust in time once the lane change commences. This is consistent with the finding of Xue et al.~\cite{C02} that rear-vehicle risk loading plays a substantial role in lane changes, and aligns with the logic of fault tree methods wherein rear-vehicle interaction is a critical node for system failure~\cite{C03,C08}.

  \item \textbf{Minimum following distance (Following\_dist\_min) consistently ranks in the Top-5 for both models.} The physical interpretation of this feature is the most direct: the closest longitudinal spacing from the preceding vehicle during a lane change intuitively reflects the urgency of longitudinal collision.
\end{enumerate}

\begin{figure}[H]
\centering
\begin{subfigure}{0.48\textwidth}
\centering\includegraphics[width=\textwidth]{fig_shap_xgb_beeswarm.png}
\caption{XGBoost}
\end{subfigure}
\hfill
\begin{subfigure}{0.48\textwidth}
\centering\includegraphics[width=\textwidth]{fig_shap_rf_beeswarm.png}
\caption{Random Forest}
\end{subfigure}
\caption{SHAP beeswarm plots for XGBoost and Random Forest.}
\label{fig:shap_beeswarm}
\end{figure}

From Fig.~\ref{fig:shap_beeswarm}, high-risk samples in both models exhibit consistent negative SHAP value distributions along the Time\_Headway\_min dimension (small headway values shifting rightward), and the distribution patterns of B\_Dist-series features are highly similar across both models, confirming the cross-model robustness of the risk factor ranking.

However, global ranking consistency does not imply sample-level feature attribution consistency. Figure~\ref{fig:shap_waterfall} presents the SHAP waterfall decomposition of XGBoost and Random Forest for the same high-risk driver sample. In the figure, $E[f(X)]$ is the baseline prediction (global mean), $f(x)$ is the final prediction for this sample, red bars indicate positive risk contributions from features, and the gray numbers on the left are the original feature values for this sample.

As shown in Fig.~\ref{fig:shap_waterfall}, both models classify this sample as high risk: XGBoost predicts a logit score of 4.413, far exceeding the baseline of $-$0.032; RF predicts a risk probability of 0.92, far exceeding the baseline of 0.332. All features contribute positively, confirming that this sample represents a persistently dangerous driving state. Notably, the core risk-driving variables differ fundamentally between the two models: XGBoost relies primarily on the TTC collision-time series (TTC\_max, TTC\_std) as the top discriminants, whereas Random Forest identifies the Time\_Headway series as the most critical risk indicators. This divergence reflects XGBoost's greater sensitivity to feature value fluctuations and extrema, making it adept at identifying transient hazardous behaviors associated with insufficient collision warning time, while RF focuses more on sustained car-following state features indicative of steady-state risk. A single tree model exhibits feature-preference limitations; fusing SHAP explanations from both models yields more complete coverage of two typical hazardous driving patterns---``sustained close-following'' and ``transient collision-time inadequacy.''

\begin{figure}[H]
\centering
\begin{subfigure}{0.48\textwidth}
\centering\includegraphics[width=\textwidth]{shap_all_exp_XGBoost_waterfall_HighRisk.png}
\caption{XGBoost (logit score)}
\end{subfigure}
\hfill
\begin{subfigure}{0.48\textwidth}
\centering\includegraphics[width=\textwidth]{shap_all_exp_RandomForest_waterfall_HighRisk.png}
\caption{Random Forest (risk probability)}
\end{subfigure}
\caption{SHAP waterfall decomposition comparison for the same high-risk sample. $E[f(X)]$: baseline prediction; $f(x)$: final prediction for this sample; red bars: positive SHAP contributions; left-side numbers: original feature values. XGBoost prioritizes TTC-series features as core risk variables, whereas RF uses Time\_Headway-series features as the primary driving factors, reflecting the two models' distinct preferences for risk feature types.}
\label{fig:shap_waterfall}
\end{figure}

\subsubsection{Cross-Location Feature Stability Analysis}

Beyond the global ranking, differential SHAP ranking analysis within each cross-location fold yields insights for location-adaptive modeling:

\begin{itemize}
  \item \textbf{Universal factors}: Time\_Headway-series features rank in the Top-3 across all five folds, with high inter-location consistency in importance ranking. This indicates that regardless of road type, the temporal spacing between the subject vehicle and the preceding vehicle remains the most fundamental lane change risk discriminant and should be incorporated into universal risk assessment models for all locations.

  \item \textbf{Conditional factors}: B\_Dist-series features rank stably in the Top-5 in Loc1--4 (subgrade and bridge segments, 120~km/h limit) but decline in ranking in Loc5 (diverging/merging area, 100~km/h limit). This shift can be explained as follows: in diverging/merging areas, the mainline traffic is compressed by merging vehicles, making the relative position of following vehicles more variable and unpredictable, thereby weakening the discriminative power of a single mean rear-distance measure. Correspondingly, speed-related features gain relative importance in this scenario.

  \item \textbf{Model consistency verification}: The high overlap between XGBoost and RF Top-5 features (intersection ratio: 4/5) confirms that the risk factor ranking derived from two independent feature attribution mechanisms is robust.
\end{itemize}

% ── 4 Discussion ──
\section{Discussion}

\subsection{Mechanisms of Cross-Location Generalization}

The substantial superiority of tree ensemble models over neural networks in cross-location generalization stems fundamentally from differences in model inductive bias. The axis-aligned partitioning of tree models confers piecewise-constant robustness under feature distribution shift: as long as test samples fall within the value range of training-set leaf nodes, reasonable classification is maintained. XGBoost's F1 degrades by approximately 2.9 percentage points from random split to cross-location evaluation (0.793 $\to$ 0.764), while RF degrades by 3.1 percentage points (0.763 $\to$ 0.732), indicating that the overall generalization loss of tree ensemble models is manageable. In contrast, MLP exhibits nearly identical random-split and cross-location F1 (0.659 vs.\ 0.660)---but this does not indicate strong generalization; rather, it reflects that MLP is globally underfitting, as its shallow fully connected architecture struggles to extract sufficient discriminative patterns from the 60-dimensional aggregate features across locations. LSTM's disadvantage arises from a dual constraint: the sample size is insufficient for adequate training of deep recurrent networks, and the cross-location distribution shift of raw temporal inputs far exceeds that of aggregate features---the latter, being hand-engineered with physically meaningful location-invariant properties (e.g., the monotonic relationship ``larger mean implies lower risk'' is preserved), whereas temporal dynamic patterns of lane changes exhibit systematic inter-location differences (e.g., acceleration buildup is shorter in Loc5 than in Loc1; drivers in Loc3 initiate lateral deviation earlier).

The contrast between random split and cross-location evaluation decomposes the generalization loss into ``model-type effects'' and ``scenario effects.'' The PCA experiment provides negative evidence: after decorrelation, LSTM F1 dropped by 9 percentage points while the standard deviation tripled, confirming that temporally correlated patterns carry critical discriminative information and that dimensionality reduction amplifies inter-location performance gaps. The enhancement experiments further demonstrate that SMOTE-synthesized samples cannot characterize the risk features of unseen locations, and that Optuna tuning gains are substantially curtailed by the heterogeneity of the diverging/merging area, indicating that diversity in road geometric configuration is the primary constraint on cross-location generalization.

\subsection{Implications for Practice and Limitations}

The separation of ``universal factors (Time\_Headway)'' from ``conditional factors (B\_Dist)'' revealed by SHAP analysis suggests a two-stage strategy for engineering deployment: ``coarse universal classification + scenario-specific fine-tuning.'' We recommend XGBoost as the base architecture, with Time\_Headway\_min serving as the universal core feature across locations, while deploying dedicated models for diverging/merging areas and enabling rapid adaptation to new locations via incremental learning.

This study has the following limitations: (1) all five locations are situated in Chongqing, where driving culture exhibits regional homogeneity; geographic generalizability requires further validation; (2) risk labels are based on SSM index-based scoring and have not been calibrated against real-world crash records; (3) 1,574 samples are modest for LSTM, and aggregate features discard precise temporal information; (4) driver heterogeneity and weather effects were not considered.

% ── 5 Conclusion ──
\section{Conclusion}

Based on 1,574 lane change samples from five scenarios of the CQSkyEyeX dataset, we systematically evaluated the cross-location generalization performance of four models---XGBoost, RF, MLP, and LSTM---for three-class lane change risk prediction.

(1) Tree ensemble models significantly outperform neural network models. XGBoost achieves optimal performance (Weighted F1 = 0.764 $\pm$ 0.079), with a cross-location degradation of approximately 2.9 percentage points, demonstrating practical feasibility for cross-location deployment. MLP shows consistently modest performance with negligible cross-location gap (0.660 vs.\ 0.659), reflecting global underfitting rather than strong generalization.

(2) Road geometric configuration is a critical moderating variable. Inter-subgrade transferability is favorable (F1 = 0.797--0.875), while generalization to the diverging/merging area drops substantially (F1 = 0.630), necessitating dedicated models for such locations.

(3) Time\_Headway\_min is the most stable cross-location risk discriminant, with a high-risk SHAP contribution (+1.251) far exceeding that of other features. The importance of B\_Dist-series features exhibits location-conditional dependence.

Future work may extend to more diverse geographic regions, explore domain adaptation and few-shot learning strategies, and incorporate driver heterogeneity and environmental factors to further narrow the cross-location generalization gap.

% ── Appendix ──
\appendix
\section{Experimental Configuration}
XGBoost (ST-optimized): 600 base learners, max depth 3, learning rate 0.095, column sampling rate 0.832, early stopping rounds 20, multi-class softmax objective. RF (ST-optimized): 600 decision trees, max depth 16, min samples per leaf 9, class\_weight = `balanced'. MLP (ST-optimized): single hidden layer with 128 neurons, ReLU activation, adaptive learning rate, early stopping regularization, StandardScaler normalization. LSTM: two recurrent layers (64 $\to$ 32 units), 25\% dropout, Adam optimizer (initial learning rate $5 \times 10^{-4}$), early stopping patience 15.

% ── Acknowledgments ──
\section*{Acknowledgments}
The CQSkyEyeX dataset used in this study was publicly released by the research group of Prof.~Jin Xu, School of Traffic \& Transportation, Chongqing Jiaotong University~\cite{C17}, to whom we express our sincere gratitude.\par
This work was supported by the National Natural Science Foundation of China (Grant No.~52372366).

% ── References ──
\bibliographystyle{elsarticle-num}
\bibliography{references}

\end{document}
